\documentclass[11pt]{article}
\usepackage[margin=1in]{geometry}
\usepackage{booktabs}
\usepackage{amsmath,amssymb}
\usepackage{microtype}
\usepackage{graphicx}
\usepackage{float}
\usepackage[hidelinks]{hyperref}
\usepackage{enumitem}
\setlist{itemsep=2pt,topsep=3pt}
\newcommand{\code}[1]{\texttt{#1}}
\newcommand{\tstar}{t^{*}}

\title{Windowed ANF Delivery:\\ Packing the DB-Mixed Circuit into ANF Blocks}
\author{local\_mixing project note}
\date{October 7, 2026}

\begin{document}
\maketitle

\noindent\emph{Status: prototype-validated pipeline stage (python tools in
\code{runs/ncdb\_ab\_20261001/}), production artifacts at $n=64$ and $n=128$
(\code{waist/final64/}, \code{waist/final128/}); Rust port planned per
\code{docs/windowed\_delivery\_design.md} \S7. All measurements cited here are
recorded in that run's \code{RESULTS.md} and in
\code{runs/exposure\_report\_20261003/REPORT.md}.}

\section{The current process}

The delivery stage consumes the \textbf{stage-3 (DB-mixed) circuit} and
produces the shipped artifact. It replaces the former stages 4--6 (split,
crossing walk, fcompress/packing). Three steps.

\subsection{Commuting reorder + window scheduling (\code{window\_pack3.py})}

The gate order of the mixed circuit is re-scheduled within its exact
commutation freedom, and simultaneously partitioned into contiguous
\textbf{windows}. Scheduling and partitioning are one optimization because
the reorder is where all the leverage is (\S\ref{sec:leverage}).

\paragraph{Commutation model (MPMCT).} Same-target writers commute; a reader
pins against \emph{all} later writers of the wire it reads. Pairwise RAW/WAR
edges are unsound for this gate set; the scheduler therefore uses epoch
conditions: per wire, a write counter; a reader of wire $c$ is ready iff
\code{writes\_done[c]} equals its required epoch, a writer iff all
earlier-epoch readers of its target have been emitted. Function preservation
is by construction and is additionally verified (256 bit-sliced input
samples, all wires) on the emitted order.

\paragraph{Value-aware planning.} Before scheduling, the stage-2 reference
circuit $A$ and the mixed circuit $B$ are simulated on 4{,}096 shared random
inputs; every intermediate value is fingerprinted. This yields the
\textbf{carried values} of $B$ (values of $A$ that re-appear on $B$'s wires)
and their \textbf{spans} (birth write $\to$ next write of the carrier wire).
A value is \emph{hidden} in the delivered artifact iff \textbf{none} of its
spans contains a window boundary; hiding is therefore all-or-nothing per
value, over all of its spans (typically ${\sim}3$, up to several hundred).

\paragraph{How the blocks are determined.} The partition is \emph{not} a
greedy accumulation of gates; greedy filling is only the background
activity. The window structure is driven by \textbf{interval reservations
and atomic bursts}, in five steps:

\begin{enumerate}
\item \textbf{Interval extraction (per span).} For each span (birth write
  $b \to$ killing write $e$) of each carried value, compute the span's
  \emph{two-sided tainted closure}: a forward taint sweep from $b$ (a gate
  joins if it reads a tainted-written wire or writes a tainted-read wire
  --- the sound any-tainted rule, since pairwise RAW/WAR edges are unsound
  for MPMCT), intersected with a backward taint sweep from $e$. The result
  is the minimal gate set that must share a window for that span to be
  hidden --- its \emph{members}. A span whose member set exceeds
  \code{W\_cap} gates disqualifies (one rescue is attempted: extend to the
  next-next writer of the carrier wire). A value is \textbf{winnable} iff
  \emph{all} of its spans have member sets $\leq$ \code{W\_cap}; values
  with any raw span longer than \code{span\_cap} are disqualified outright.
\item \textbf{Readiness bookkeeping (epochs).} Every gate carries an unmet
  count: a reader of wire $c$ becomes ready when $c$ reaches the
  write-epoch the original order gave it; a writer of wire $t$ when all
  earlier readers of $t$ are emitted. This is the exact commutation
  freedom --- the scheduler can emit any ready gate at any time without
  changing the function.
\item \textbf{Drain phase (the greedy part).} Intervals are processed by
  ascending end position. Between bursts the scheduler emits
  \textbf{filler}: ready gates that belong to \emph{no} living interval
  (protection count zero). Filler fills windows with material whose
  placement is exposure-irrelevant.
\item \textbf{Atomic bursts + value-level accounting.} When an interval's
  turn comes, the window is first checked: if the remaining gate capacity
  is below the interval's size ($+$ a pad of 4) or the composed ANF is
  above $\tfrac34$ of the monomial budget, the window is closed
  \emph{pre-emptively} so the burst is never split by a scheduled close.
  Then \textbf{all member gates of the interval are scheduled
  consecutively} (dependencies pulled from the ready heap as needed)
  inside the one window. Hiding is all-or-nothing per value: the first
  touched span commits the value; if any span dies --- a window closed
  mid-burst, or a member was consumed earlier as filler/another burst
  (\code{resumed\_later}) --- \textbf{all sibling intervals of that value
  are killed} (\code{kill\_value}) and their reserved gates are released
  to the filler pool immediately, so no capacity is spent on half-hidden
  values. After each burst, \textbf{contagion} completes any committed
  interval with $\leq 24$ remaining members in place while the window is
  still open.
\item \textbf{Budget-adaptive close.} The close condition measures the
  \emph{real} deliverable cost: the window's sparse ANF is composed
  incrementally gate-by-gate during scheduling, and the window closes when
  the monomial count exceeds the soft budget \code{B\_mono} (hard
  $2\cdot\code{B\_mono}$ inside a burst) or the gate count exceeds the
  soft target $G$ (hard $2G$ inside a burst). Every interval death is
  attributed (\code{drain\_close} / \code{burst\_close} /
  \code{resumed\_later} / \code{sibling}), which is how scheduler losses
  are diagnosed (\S2.4).
\end{enumerate}

After the main pass, remaining gates are emitted by readiness, the final
window closes, the emitted order is re-verified functionally, and the ANF
emission stage (\S1.2) re-composes each window independently against the
larger emission budget with per-window verification.

\subsection{ANF emission (\code{anf\_compose.py})}

Each window is delivered as the canonical algebraic normal form of its
touched wires: for every wire written inside the window, the exact
$\mathrm{GF}(2)$ polynomial of its value at the window's end as a function of
the window's \emph{input} wire values. Composition is sparse and exact, with
true $\mathrm{GF}(2)$ cancellation (per-monomial toggle ---
symmetric-difference dedup is a measured bug, caught by re-verification).
Monomials and wires are canonically ordered; the within-window gate text,
gate count, gate order, and all intermediate spellings are discarded.

Budget: 16{,}384 monomials per window. Measured usage is far below it: median
184 ($n{=}64$) / 207 ($n{=}128$) monomials per window, p90 $\approx$
290--430, maximum 7{,}673, \textbf{zero overflows} at both scales. Every
window is functionally re-verified against the gate-level window on the
shared sample set before serialization.

\subsection{Serialization}

The artifact (\code{deliverable.bin}) is the sequence of window ANFs in a
binary format (wire ids, monomial lists over window inputs). Nothing else
ships: no gate list, no boundary metadata beyond the windows themselves, no
stage provenance.

\begin{table}[H]\centering
\begin{tabular}{lll}
\toprule
 & $n=64$ ($|C|=1{,}000$) & $n=128$ ($|C|=3{,}000$) \\
\midrule
stage-2 circuit & 771{,}029 g / 256 w & 2{,}240{,}757 g / 512 w \\
stage-3 (mixed) input & 1{,}648{,}679 g & 4{,}803{,}546 g \\
windows & 10{,}702 & 28{,}986 \\
ANF payload & 24.6 MB ($\times$3.88 vs gate list) & 93.7 MB ($\times$5.06) \\
\code{deliverable.bin} & 27.5 MB & 103.1 MB \\
xz-compressed & 4.52 MB & 15.3 MB \\
\bottomrule
\end{tabular}
\end{table}

Evaluation cost (python reference evaluator, 1{,}024-lane bit-sliced):
0.68 / 1.70 ms per input in gate form, 0.85 / 2.54 ms in ANF form, bit-exact
agreement between the two forms at both scales. Compressibility is
6.1--6.7$\times$ under xz --- and the compressed ANF is \emph{smaller} than
the compressed gate list, i.e.\ the $\times$4--5 raw inflation is redundancy
of the format, not information added.

\subsection{Full production parameters, both scales}

Stage-2 knobs are identical at both scales (the managed embedded-masking
preset): $K{=}2$ band wires per LGI, \code{max\_open} 3 / \code{min\_open}
2, quad-fire, balanced masks and balanced seed, rich band seed
(\code{seed\_fill\_nl} 4), feistel slice guards (all gates $\leq 2$
controls), carrier shuffling dose 8 with \textbf{persistent migration}
(\code{EMBEDDED\_MASKING\_SHUFFLING=8}, \code{...\_RETURN\_HOME=0}).
Stage-3 flags are identical: \code{--tdp --db-mixing --profile 2,2,2,2,2
--p-db 1.0 --p-comp 1.0 --p-any 0.1 --db-max-degree 9 --db-max-span 30
--db-wire-terms 1024 --db-total-terms 2048}, regular frozen store
\code{frozen\_m1\_m11} (\code{--p-comp}/\code{--p-any} are retired since
the four-stage cleanup: they were inert at \code{--p-db 1.0}, and current
builds reject them), moves $= \lceil 2\cdot 2\cdot 1.3\cdot
(\text{stage-2 gates})\rceil$. Packing and emission are identical:
\code{W\_cap} 128, $G$ 128 (hard 256), \code{B\_mono} 6{,}144 (hard
12{,}288 in a burst; $\tfrac34$ pre-burst guard), \code{span\_cap}
3{,}000, contagion threshold 24, pad 4; emission budget 16{,}384
monomials/window; 4{,}096 fingerprint lanes. Fresh CSPRNG seeds per run,
recorded in each run directory.

\begin{table}[H]\centering
\small
\begin{tabular}{lll}
\toprule
parameter & $n=64$ run & $n=128$ run \\
\midrule
source $C$ & 64 wires, 1{,}000 gates & 128 wires, 3{,}000 gates \\
sandwich ($|D|$, $s$, slice gates) & 1{,}000 / 384 / 1{,}280 & 3{,}000 / 896 / 2{,}560 \\
sandwich circuit & 2{,}832 g / 128 w & 7{,}920 g / 256 w \\
stage-2 output & 771{,}029 g / 256 w & 2{,}240{,}757 g / 512 w \\
mixer moves & 4{,}009{,}351 & 11{,}651{,}936 \\
stage-3 output (wall time) & 1{,}648{,}679 g ($\sim$26 min) & 4{,}803{,}546 g ($\sim$76 min) \\
carried / winnable & 145{,}481 / 131{,}226 (90.4\%) & 426{,}891 / 382{,}071 (89.5\%) \\
hidden by construction & 125{,}592 (95.7\% of winnable) & 337{,}934 (88.4\%) \\
windows & 10{,}702 & 28{,}986 \\
window monomials med/p90/max & 194 / 301 / 7{,}682 & 207 / 332 / 2{,}846 \\
exposed distinct at boundaries & 23{,}347 (16.0\% of carried) & 89{,}072 (20.8\%) \\
\bottomrule
\end{tabular}
\end{table}


\section{Rationale, tradeoffs, and what the process optimizes}

\subsection{The threat model that motivates windowing}

Every exposure metric in the 2026-10 program --- per-value visibility, the
$\tstar$ exposure census (exact affine predictors), the equality-cluster
attack --- quantifies over \textbf{cuts of the delivered spelling}: the
adversary reads wire values at positions between delivered blocks. Two
measured laws made gate-granular delivery look fundamentally leaky:

\begin{itemize}
\item \textbf{Reader law}: any value with $\geq 1$ bare reader is
  ${\sim}100\%$ visible under every-gate cuts (measured
  $3.2\,/\,55.4\,/\,63.8\,/\,100.0\%$ visibility at fanout
  $0/1/2/{\geq}3$ on the TDP, stage 3).
\item \textbf{Closure law}: function preservation forces interior basis
  changes to close (the second bracket undoes the first), so \textbf{no
  self-closing interior move can raise the min-over-cuts predictor weight at
  any dose} --- a CNOT/transvection twist sweep that bloated the circuit
  $\times 37$ left $\tstar{=}1$ mass at 84\%. The interior move system
  contains no weight-raising operator; the $\tstar$ spectrum belongs to the
  circuit ends.
\end{itemize}

Windowing attacks the quantifier instead of the values: \textbf{a window has
no interior cuts.} A value is exposed only if it is alive across a window
boundary. Deleting cuts (a) makes read values hideable for the first time
--- the reader law presumed a cut between every pair of gates; (b) sidesteps
the closure law --- nothing is rewritten between cuts, the cuts are deleted;
and (c) annihilates all within-window text: templates, $k$-grams, gate
provenance, stage-2 vocabulary signatures, and the gate count itself vanish
into a canonical polynomial. What remains observable is exactly: values alive
at boundaries, the boundary equality-linkage graph, and the outputs.

\subsection{The objective function}

The packer maximizes the number of \textbf{distinct carried stage-2 values
hidden} --- ground truth measured by fingerprint simulation on the emitted
order, not a proxy. Secondary structure it respects:

\begin{itemize}
\item value-level all-or-nothing (a value with one exposed span is exposed;
  partial hiding is worthless);
\item the C-view invariant: the source computation's interior must stay at
  \textbf{zero} exposure (measured: the output-only floor --- exactly the 64
  / 128 outputs visible, all $\tstar{=}1$ --- holds through the entire
  delivery stage at both scales);
\item ANF budget and window caps (digestibility and size).
\end{itemize}

\subsection{Where the leverage is: reorder, not boundary placement}
\label{sec:leverage}

Measured on the TDP: delivered-order carrier spans are long and waists are
flat --- optimal boundary placement in the \emph{given} order is
approximately uniform placement, worth almost nothing. But a value's minimal
bare residency under commuting reorder is its \textbf{DAG interval}, median 6
gates on stage-3 material --- $\times 21$ \textbf{shorter} than its spelled
span. The entire mechanism is: reorder to compress each value's bare
residency to its DAG interval, then drop windows around the intervals.

\subsection{The window-size tradeoff}

The one real tradeoff surface is window size $G$:

\begin{itemize}
\item \textbf{Hiding rises ${\sim}{+}14$ points per doubling of $G$}
  (measured sequence, migrated material: 46\% hidden @$G{=}64$, 58\%
  @$G{=}128$, 72\% @$G{=}256$ --- 5.2\% of the full stage-2 vocabulary still
  exposed at $G{=}256$).
\item The \textbf{ANF budget never binds}: median monomial usage is
  ${\sim}200$--$400$ of multi-thousand budgets even at 512-gate windows,
  zero overflows anywhere in the program. The limit on $G$ is architectural,
  not combinatorial: boundary count vs the circuit ends, the
  re-synthesis-hardness floor (very small windows are canonically
  re-synthesizable; very large windows reduce the boundary count toward the
  ends), and delivery size ($\approx\times 12$ at $G{=}512$).
\item Scheduler quality, not material, is the current binding constraint: on
  migration-refreshed material, 93\% of carried values are winnable while v3
  hides 72--73\% of winnable; the gap is attributed by the death counters to
  burst-time closes and sibling cascades (scheduler v4, value-unit planning
  with conflict-graph pre-grouping, is the specified next step; estimated
  ceiling 94\%).
\end{itemize}

Packer evolution, same ground truth (\% of carried values hidden, stage-3
input, $W{=}20$ era $\to$ $G$-era): naive uniform windows 5\% $\to$ v1
interval-burst scheduler 21\% $\to$ v2 value-level all-or-nothing 32\% $\to$
v2.1 35\% $\to$ v3 budget-adaptive 46/58/72\% at $G=64/128/256$.

\subsection{Costs accepted}

\begin{itemize}
\item \textbf{Size} $\times$4--5 raw ($\times{\approx}0.9$ after
  compression, \S1.3): the format buys erasure, not compression.
\item \textbf{Boundaries are new visible structure}: the window lattice
  itself is observable. Measured consequence: a value recurring $\geq 5$
  times among the carried values is \emph{always} boundary-exposed
  (deterministic at both scales; non-exposed maximum multiplicity is 4).
  This ``recurrence law'' shifted the program's focus upstream to stage-2
  carrier hygiene (the celebrity program: migration, hot-targeting,
  threshold retirement), since no boundary placement can hide a value whose
  spans tile the stream (celebrity liveness was measured at ${\sim}110$ of
  256 wires at \emph{every} cut position).
\item \textbf{Affine-metric caveat}: all exposure numbers quantify exact
  affine predictors at boundaries; the approximate/nonlinear-predictor flank
  is an open thread of the exposure program (known empty for the QC-style
  detector: agreement $\geq 0.98$ coincides with exact equality).
\end{itemize}

\section{Why stages 4--6 were abandoned}

The decision rests on the stage-wise exposure study (nine arms $\times$ four
stages, random family and TDP, 2026-10-03) plus the direct input-choice
measurement of the windowing program (2026-10-04).

\subsection{What stages 4--6 actually bought (measured)}

Per-stage visibility \% [$\tstar{\geq}2$ share \%], random-family arms,
every-cut grids (excerpt; full table in
\code{runs/exposure\_report\_20261003/REPORT.md}):

\begin{table}[H]\centering
\begin{tabular}{lcccc}
\toprule
arm & s3 mix & s4 split & s5 crossing & s6 fcompress \\
\midrule
random $f{=}2$ & 85.1 [1.5] & 84.0 [8.3] & 76.1 [8.9] & 79.2 [5.8] \\
$f{=}10$ & 79.5 [1.7] & 79.0 [5.1] & 76.0 [6.1] & 76.3 [4.3] \\
eq.\ arity-2 & 82.6 [1.3] & 81.6 [4.5] & 78.7 [6.0] & 79.4 [4.6] \\
twists $p{=}.05$ & 87.0 [2.2] & 87.0 [5.8] & 82.7 [7.8] & 83.9 [5.8] \\
\bottomrule
\end{tabular}
\end{table}

Stage signatures, replicated at TDP scale:

\begin{itemize}
\item \textbf{Stage 4 (split)} is a visibility-neutral $t$-raiser: it moves
  4--8 points of $\tstar{=}1$ mass into small parities and erases nothing.
  In carrier terms it \emph{adds} exposure surface: 105.7k $\to$ 107.9k
  carried values on the TDP.
\item \textbf{Stage 5 (crossing)} was the only post-processing eraser:
  $-3\ldots{-}5$ points of visibility in every arm (105.7k $\to$ 89.8k
  carried on the TDP). This is the entire measured security content of the
  old post-3 pipeline.
\item \textbf{Stage 6 (fcompress)} partially \emph{refunds} stages 4--5:
  $+1\ldots{+}3$ points, parities collapse back. Its packing component is
  exposure-neutral (pack-only $\equiv$ crossing to $\pm 0.1$;
  pack-then-compress $\equiv$ compress-then-pack, order flip null at $+2\%$
  size).
\item At TDP scale the whole chain is flat: stage-2-view visibility
  $26.5 \to 26.5 \to 25.6 \to 27.2/27.3\%$ across s3 $\to$ s6. The C-view
  output-only floor is owned by stage 2 and never moves.
\end{itemize}

So the old pipeline's net effect was: one eraser (crossing), one neutral
$t$-raiser (split) whose gains the compressor refunded, at the price of
large inflation and run time.

\subsection{The direct comparison: windowing s3 vs windowing s5}

The windowing program's first decision was its input. Measured:

\begin{itemize}
\item \textbf{Crossing welds the dependency structure.} Its $\times 3.45$
  inflation stretches the median DAG interval from 6 gates (s3) to 26 (s5)
  --- $\times 4$ --- so the reorder lever that windowing lives on is largely
  destroyed. Per-value hideable ceilings at matched window size $W{=}20$:
  69.6\% of carried values (s3) vs 46.9\% (s5).
\item \textbf{Crossing's erasure does not pay for the welding.} In absolute
  terms at $W{=}20$: ${\sim}32$k exposed values on the s3 ceiling vs
  ${\sim}48$k on the s5 ceiling --- windowing the raw mixed circuit strictly
  dominates windowing the crossed one, despite crossing's 15\% carrier
  erasure.
\item \textbf{Split only adds carriers} (\S3.1), and the walk family
  generally is exposure-hostile: the order-preserving greedy walk measured
  97.4\% exposure against its own 79.4\% floor --- verbatim-dressing passes
  are exposure-\emph{maximal}.
\end{itemize}

Against these, a single windowing pass hides 58--72\% of carried values
($G{=}128$--$256$) --- an order of magnitude beyond crossing's
$-3\ldots{-}5$ points --- while \emph{deleting} within-window text rather
than respelling it, at lower wall-clock cost than the crossing walk it
replaces.

\subsection{Division of labor after the replacement}

\begin{itemize}
\item \textbf{Stage 2} --- semantic hiding: which values ever exist bare
  (data path at zero bare reads; band carrier hygiene: lifetime, migration,
  dedup/retirement, seed density).
\item \textbf{Stage 3} --- un-invertible respelling, dilution, DB-form
  digestion.
\item \textbf{Windowed ANF delivery (new stage 4)} --- cut deletion:
  \emph{where} anything is observable at all, plus canonical within-window
  respelling.
\item \textbf{Circuit ends (sandwich/slicing)} --- the $\tstar$-spectrum of
  the survivors (interior moves provably cannot move it).
\end{itemize}

\section{Production reference results (pipeline v2.1, final64/final128)}

End-to-end: classic sandwich $\to$ embedded masking (rich band seed,
persistent carrier migration dose 8, read-pressure partners, feistel slice
guards) $\to$ db\_mixing $\to$ window-pack $G{=}128$ $\to$ per-window sparse
ANF. Function checks green at every stage, both scales; all windows
re-verified; zero budget overflows; fresh CSPRNG seeds recorded in the run
directories.

\begin{table}[H]\centering
\footnotesize
\begin{tabular}{p{4.6cm}p{5.2cm}p{4.2cm}}
\toprule
leakage parameter & $n=64$ & $n=128$ \\
\midrule
C-view (the secret) & 6.4\% = exactly the 64 outputs, $\tstar{=}1$ &
  4.3\% = the 128 outputs ($+2$ interior $\tstar{=}4$) \\
stage-2-view visible (exact lattice) & 27.3\% ($\tstar{=}1$: 22.1,
  $\tstar{\geq}2$: 5.2) & 28.4\% (23.6 / 4.7) \\
distinct values exposed at boundaries & 23{,}347 = 16.0\% of carried =
  3.8\% of vocab.\ & 89{,}072 = 20.8\% / 5.0\% \\
class uniformity (data/mid/band) & 26.5 / 27.3 / 27.7\% & 28.3 / 27.2 /
  28.7\% \\
non-exposed max multiplicity & 4 & 4 \\
max boundary equality cluster (proxy) & 342 & 397 \\
\bottomrule
\end{tabular}
\end{table}

The residual exposed mass is dominated by high-recurrence band vocabulary
(the celebrity class, 330/966 values of multiplicity $\geq 257$), which is a
\textbf{stage-2 birth phenomenon}, not a packing phenomenon: boundary
redrawing was measured infeasible (celebrity spans median 494 / p95 2{,}505
packed gates, liveness ${\sim}110$ wires at every cut position), and the fix
in progress (threshold retirement of band carriers at the emitter) operates
upstream of this stage. The delivery stage's own contract --- C-view at the
output-only floor, within-window annihilation, recurrence-law behavior at
boundaries --- is identical with and without that fix.

\section{Artifacts and tools}

\begin{itemize}
\item Design doc: \code{docs/windowed\_delivery\_design.md} (architecture +
  the stage-2 band-subsystem spec).
\item Tools: \code{runs/ncdb\_ab\_20261001/window\_pack3.py},
  \code{anf\_compose.py}, \code{boundary\_cluster.py},
  \code{waist/recurrence2.py}, \code{waist/anf\_dump\_eval.py},
  \code{wire\_balance.py}; exposure scans via \code{hmap\_affine --exposure
  --g-cuts} (boundary-lattice mode).
\item Production artifacts:
  \code{runs/ncdb\_ab\_20261001/waist/\allowbreak final64/} and
  \code{final128/} (stage-2/3 circuits, packed order + bounds,
  deliverable.bin + compressed forms, full census logs, seeds).
\item Measurement record: \code{runs/ncdb\_ab\_20261001/}\allowbreak\code{RESULTS.md};
  stage-wise study:\\ \code{runs/exposure\_report\_20261003/REPORT.md}.
\end{itemize}

\end{document}
